% 第 17 题频率分布直方图：柱高按固定源 PDF 像素实测重绘（组距 0.5，共 9 组）。
% 最左组 [0,0.5) 频率为 0（原图如此）：纵轴与第一根柱之间留一个组距空档，不画柱。
\begin{tikzpicture}[>=stealth,every node/.style={font=\small}]
\def\ux{1.15}   % 横轴：1 立方米
\def\uy{3.6}    % 纵轴：单位 1（频率/组距）
% 坐标轴（横轴标题置于轴箭头右端之后，与原图一致；不可用 below left 锚到 (5.45,0)——会覆盖 2.5–4.5 刻度）
\draw[->] (0,0)--(5.45,0) node[right,yshift=-2pt] {用水量（立方米）};
\draw[->] (0,0)--(0,2.05);
\node[anchor=south west] at (0.10,2.02) {\begin{tabular}{c}频率\\\hline 组距\end{tabular}};
% 刻度虚线（0.1–0.5）与纵轴刻度标签
\foreach \v in {0.1,0.2,0.3,0.4,0.5} {
  \draw[dashed] (0,\v*\uy)--(5.1,\v*\uy);
  \node[left] at (-0.08,\v*\uy) {\(\v\)};
}
% 横轴刻度标签
\foreach \t in {0,0.5,1,1.5,2,2.5,3,3.5,4,4.5} {
  \node[below] at (\t*\ux,-0.10) {\(\t\)};
}
% 柱（空心矩形，相邻柱共用竖边；最左组无柱）：
% [0.5,1)=0.2，[1,1.5)=0.3，[1.5,2)=0.4，[2,2.5)=0.5，[2.5,3)=0.3，[3,3.5)=[3.5,4)=[4,4.5)=0.1
\foreach \i/\h in {1/0.2, 2/0.3, 3/0.4, 4/0.5, 5/0.3, 6/0.1, 7/0.1, 8/0.1} {
  \draw ({0.5*\i*\ux},0) rectangle ({0.5*(\i+1)*\ux},{\h*\uy});
}
\end{tikzpicture}
